% This file should be rename with the speaker's last name. (e.g : smith)

% Write the author names in the order you wish them to appear
% and underline the one who will give the talk.

% Only the speaker's email address is required.

\begin{abstract_online}{Some examples of calculation improper integrals using CAS}{%
        \underline{W{\l}odzimierz Wojas}$^{1}$, \underline{Jan Krupa}$^{1}$}{[jan\_krupa@sggw.pl]
        }{%
        $^1$ Department of Applied Mathematics, Warsaw University of Life Sciences (SGGW), Warsaw, Poland}


\newcommand{\ds}{\displaystyle}
\newcommand{\dx}{\,\mathrm{d}}


Improper integrals are taught students in framework of such academic courses as: Calculus, Mathematical Analysis or Higher Mathematics as a standard. In this talk we would like to present some didactics examples representing different approaches to calculate improper integrals using Mathematica and wxMaxima. We will present two examples of improper integral calculated using Riemann sums. We will compare Riemann and Lebesgue approaches to integral  $\ds \int_0^{\infty}\frac{\sin x}{x}\dx{x}$. We will also analyse complex approach to calculate improper integrals on the following examples: $\ds \int_{-\infty}^{\infty}\frac{x^2}{x^6+6x^4+9x^2+4}\dx{x}$   and $\ds \int_L\frac{\re z}{\bar{z}}\dx{z}$ where $L$  is a broken line $ABC$  on Gauss plane and $A=-1$, $B=0$, $C=i$.


\textbf{Keywords} \newline{} Higher education, Improper integrals, Application of CAS, Mathematica, Mathematical didactics




\textbf{References}\newline{}
[1] \textsc{I.M. Roussos}, \emph{Improper Riemann Integrals}, CRC Press, 2014\newline{}
[2] \textsc{Jian-Ke Lu, Shou-Guo Zhong and Shi-Qiang Liu}, \emph{Introduction to the Theory of Complex Functions}, Series of Pure Mathematics, World Scientific, 2002
\end{abstract_online}